\section{Compiler setup}
\label{section:preparation:compiler-settings}

Before we start, it is important to doublecheck whether we have used the correct
compiler settings. It makes no sense to benchmark a code which is not producing
tailored code for our particularly machinery (therefore it is important to
benchmark on the right system), and it makes no sense to benchmark a code that
is not using the latest compiler features.


First of all, we work with the the most aggressive optimisation that our
compiler can offer. The minimum setup is to compile with 

\begin{verbatim}
 -O3
\end{verbatim}


%  CFLAGS="-g" CXXFLAGS="-g" --with-multithreading=tbb_extension
% --enable-blockstructured --enable-particles --enable-loadbalancing --enable-finiteelements --enable-mghype --enable-swift --enable-exahype CXXFLAGS="-std=c++20 -O3 -g3 -DTBBExtensionAddRanges -fno-exceptions -fiopenmp -Wno-unknown-attributes" LDFLAGS="-fiopenmp" LIBS="-ltbb_debug" CXX=icpx CC=icpx

Second, we create code that is tailored towards our particular machinery.
The correct setting for this depends on 



Finally, we might want to add debug symbols.
These instruct the compiler to insert additional information into the source
code which tools can then later use to \marginpar{Rueckschluss ziehen} 
from which source code line a certain instruction stems from.
Usually, you cannot know from a single instruction in a machine code, where this
comes from in the original code.
With the debug systems, we inject exactly this information into the executable.
Therefore, we should add \texttt{-g}.
Some compilers support more detailed information (e.g.~using \texttt{-g3}).
It requires studying the compiler's feedback whether these flags disable
optimisations, in which case we might have to balance insight vs.~best-case
speed.
